\section{Related Work}
\label{sec:related}

\subsection{Deep Learning for Plant Foliar Pathology}
The automation of foliar plant disease identification has evolved markedly over the past decade. Early methodologies relied on handcrafted texture descriptors, such as Gray-Level Co-occurrence Matrices (GLCM), local binary patterns (LBP), and Scale-Invariant Feature Transforms (SIFT), paired with shallow classifiers such as Support Vector Machines (SVM) \cite{kamencay2020improved}. The landmark study by Mohanty \textit{et al.} \cite{mohanty2016using} catalyzed the transition to deep learning, training AlexNet and GoogLeNet on 54,306 laboratory images from the PlantVillage dataset to achieve 99.35\% diagnostic accuracy. Subsequent works by Ferentinos \cite{ferentinos2018deep} and Fuentes \textit{et al.} \cite{fuentes2017robust} expanded convolutional classification and object detection architectures across multi-crop greenhouse databases.

However, as rigorously documented by Barbedo \cite{barbedo2019plant}, models trained on pristine laboratory specimens experience precipitous performance declines of 20\% to 35\% when evaluated in operational orchards. Laboratory photographs typically isolate excised leaves on uniform black or white cards under shadowless lighting. In contrast, in-situ orchard images feature severe background interference from soil, weeds, and overlapping non-target foliage, alongside extreme illumination gradients caused by direct solar glare or canopy shadows. Consequently, contemporary research has pivoted toward developing architectures resilient to field-induced visual clutter.

\subsection{Lightweight and Edge Architectures for Agriculture}
Deploying vision models in rural agricultural environments necessitates strict adherence to embedded compute and power envelopes. Standard deep architectures, such as ResNet-50 \cite{he2016deep} (25.56\,M parameters, 4.12\,GFLOPs) and Vision Transformers \cite{dosovitskiy2020image, liu2021swin} ($>28\,\text{M}$ parameters), exceed the storage and real-time execution capacities of mobile diagnostic terminals and low-cost single-board computers (SBCs).

To reconcile capacity with efficiency, researchers have embraced lightweight architectures. Sandler \textit{et al.} \cite{sandler2018mobilenetv2} introduced inverted residual blocks with linear bottlenecks in MobileNetV2, significantly decreasing memory access costs. Tan and Le \cite{tan2019efficientnet} proposed compound coefficient scaling in EfficientNet, balancing network depth, width, and input resolution. Howard \textit{et al.} \cite{howard2019searching} leveraged hardware-aware Neural Architecture Search (NAS) and lightweight Squeeze-and-Excitation modules to formulate MobileNetV3. Similarly, Ma \textit{et al.} \cite{ma2018shufflenet} devised ShuffleNetV2, employing channel split and shuffle operations to minimize memory access overhead. While these architectures deliver rapid inference on edge hardware, their depthwise separable convolutions prioritize channel-wise spatial filtering over high-frequency spatial-frequency localization, rendering them susceptible to texture blur under adverse field conditions \cite{ahmed2022field}.

\subsection{Spatial-Frequency and Gabor Filter Representations}
Gabor wavelets provide mathematically optimal joint localization in both the 2D spatial and spatial-frequency domains, reaching the theoretical lower limit of the Heisenberg-Gabor uncertainty relation \cite{daugman1985uncertainty}. Because biological simple cells in the mammalian primary visual cortex exhibit response profiles closely modeled by 2D Gabor functions, Gabor filtering has historically served as a foundational technique for texture segmentation, fingerprint recognition, and face verification \cite{manjunath1996texture}.

In agricultural computer vision, classical Gabor filters have been widely employed to extract directional leaf venation and punctate lesion patterns \cite{wang2021gabor}. However, standard pipelines rely on handcrafted, static filter banks where orientations, frequencies, and Gaussian window parameters are fixed a priori. Recently, efforts have emerged to incorporate Gabor functions into neural networks. Luan \textit{et al.} \cite{luan2018gabor} proposed Gabor Convolutional Networks (GCNs), modulating convolutional filter weights with harmonic Gabor functions to bolster rotation invariance on benchmark image datasets. Alekseev and Bobe \cite{alekseev2019gabor} demonstrated that learnable Gabor parameterizations accelerate training convergence in early convolutional layers. Nevertheless, existing hybrid approaches either freeze Gabor parameters throughout optimization or enforce rigid parameter sharing across layers without dedicated cross-branch feature calibration.

\subsection{Litchi Disease Recognition and the BDLitchi Benchmark}
Despite the substantial agro-economic value of litchi, automated diagnostic research specific to this crop has remained sparse compared to staples such as rice, wheat, and maize. Early investigations utilized small private datasets (rarely exceeding 500 images) evaluated under laboratory conditions with standard ResNet or VGG backbones. 

A major advancement was achieved with the public release of the BDLitchi dataset by Islam \textit{et al.} \cite{litchi_dataset_2022}, comprising \totalImages{} natural field-condition images spanning \numClasses{} foliar conditions captured in the commercial orchards of Bangladesh. Alongside the dataset, the authors proposed LitchiChebNet, a spectral graph convolutional network utilizing Chebyshev polynomial approximations over segmented leaf boundary superpixels, achieving a reported Top-1 accuracy of 96.40\%. While innovative, LitchiChebNet entails significant drawbacks: the graph construction and superpixel segmentation stages impose high computational latency (averaging 78\,ms per image), and the method is vulnerable to topological failure when severe disease necrosis or insect chewing damages leaf margins.

\subsection{Research Gap and Synthesis}
Table~\ref{tab:related} summarizes the literature landscape, comparing prior methodologies, evaluation domains, reported accuracies, and operational limitations. 

\begin{table*}[t]
\centering
\caption{Systematic Summary of Related Literature in Foliar Crop Pathology and Gabor-Deep Hybrid Models}
\label{tab:related}
\begin{tabular}{p{3.2cm}p{3.8cm}p{3.2cm}cc}
\toprule
\textbf{Study} & \textbf{Methodological Approach} & \textbf{Benchmark Dataset} & \textbf{Accuracy (\%)} & \textbf{Primary Limitation Addressed in this Work} \\
\midrule
Mohanty \textit{et al.} (2016) \cite{mohanty2016using} & AlexNet / GoogLeNet & PlantVillage (Lab) & 99.35 & Evaluated solely on uniform laboratory backdrops; fails in orchards \\
Ferentinos (2018) \cite{ferentinos2018deep} & VGG / ResNet-50 & PlantVillage + Greenhouse & 99.53 & High parameter count (>98 MB); susceptible to natural solar shifts \\
Barbedo (2019) \cite{barbedo2019plant} & Individual Lesion CNN & Multi-Crop Field Set & 84.10 & Requires manual lesion segmentation before inference; unscalable \\
Wang \textit{et al.} (2021) \cite{wang2021gabor} & Fixed Gabor + DenseNet & Apple Leaf Benchmark & 96.80 & Handcrafted static Gabor bank cannot adapt to lesion orientation \\
LitchiChebNet (2022) \cite{litchi_dataset_2022} & Chebyshev Graph CNN & BDLitchi Raw & 96.40 & Graph building adds 78 ms latency; vulnerable to boundary occlusions \\
Luan \textit{et al.} (2018) \cite{luan2018gabor} & Gabor Convolutional Net & PASCAL VOC / CIFAR & 92.50 & Strict filter parameter tying without cross-attentive texture fusion \\
Ahmed \textit{et al.} (2022) \cite{ahmed2022field} & MobileNetV2 Edge & Rice Disease Field Set & 95.20 & Standard depthwise kernels degrade by 22\% under optical motion blur \\
\midrule
\textbf{LitchiHybridNet (Ours)} & \textbf{Dual-Branch Learnable Gabor + MobileNetV3 + GAFM} & \textbf{BDLitchi Leakage-Free} & \textbf{99.04} & \textbf{Provides robust frequency inductive bias and real-time edge speed} \\
\bottomrule
\end{tabular}
\end{table*}


The literature review exposes three major research gaps:
\begin{enumerate}
    \item \textbf{Absence of End-to-End Learnable Frequency Priors:} Existing agricultural models either discard spatial-frequency texture priors entirely or rely on static, handcrafted Gabor filters that cannot adapt to the morphological variability of natural lesions.
    \item \textbf{Suboptimal Feature Fusion:} Prior multi-stream architectures rely on simple vector concatenation or heavy multi-head self-attention, failing to strike an effective balance between discriminative cross-modal calibration and edge computational efficiency.
    \item \textbf{Pervasive Data Leakage in Benchmark Splits:} Agricultural field datasets frequently suffer from unaddressed burst-shot duplicate leakage, where near-identical consecutive frames span both training and testing sets, artificially inflating reported baseline metrics.
\end{enumerate}

LitchiHybridNet directly addresses these deficiencies through an end-to-end differentiable learnable Gabor layer, an edge-efficient cross-gated attention fusion module (GAFM), and a rigorously audited, leakage-free benchmark evaluation protocol.
